% !TeX root = ../main.tex
\section{背景与动机}
\label{sec:background}

\subsection{浮点精度与微缩放格式}
\label{sec:formats}

BF16/FP16 由逐元素的指数与有效数确定动态范围和精度~\cite{burgess2019bfloat16}；微缩放格式则将低位宽元素与共享尺度共同编码（图~\ref{fig:fpformats}）。OCP MXFP8/6/4 每 32 个元素共享一个 E8M0 尺度~\cite{ocpmx2023}；NVFP4 使用 16 元素一组的 E4M3 尺度~\cite{nvidia2024nvfp4}；HiF4 使用 64 元素一组的无符号 E6M2 基础尺度，并为每 8、4 个元素分别配置一位微指数~\cite{hifloat4}。元素编码、分组粒度和尺度编码分别决定解包、局部规约与尺度融合的硬件需求。

\begin{figure}[t]
\centering
\includegraphics[width=\columnwidth]{fp_format.png}
\caption{元素编码与共享尺度~\cite{ocpmx2023,nvidia2024nvfp4,hifloat4}。字段标注为位宽，$g$ 为缩放组大小。}
\label{fig:fpformats}
\end{figure}


\subsection{张量计算与点积单元}
\label{sec:tensor_organization}

对于矩阵乘累加 $D=C+AB$，$M$、$N$ 是输出维度，$K$ 是归约维度。点积单元先对齐多个乘积并集中归约，再统一规格化和舍入，以共享后端摊薄逐项浮点处理的成本~\cite{lutzFusedFP84Way2024}。本文采用 H 类原生归约长度为 16 的 DP 作为基线，与已有数值建模研究识别的点积累加粒度相对应~\cite{xie2025mmasim}。完整 GEMM 沿归约维拆为多笔请求，单 DP 接收显式 $C$ 并输出 $D$，跨请求部分和保存在其外部。

有无组尺度的格式可用同一局部点积结构描述。依据 OCP 的点积定义，共享尺度可提出组内归约~\cite{ocpmx2023}。将 $L$ 项乘积按缩放组边界划分为 $G=\lceil L/g\rceil$ 组，末组不足部分补零，加入累加输入 $c$ 后，有
\begin{equation}
d=c+\sum_{j=0}^{G-1}X_j^{(A)}X_j^{(B)}
\underbrace{\sum_{i=0}^{g-1}P_{j,i}^{(A)}P_{j,i}^{(B)}}_{q_j}.
\label{eq:blockdot}
\end{equation}
其中，$X_j^{(A)}$、$X_j^{(B)}$ 为两路输入第 $j$ 组的共享尺度，$P_{j,i}^{(A)}$、$P_{j,i}^{(B)}$ 为组内元素的解码值，$q_j$ 为组内点积；HiF4 的局部微指数计入元素解码。对于无组尺度的格式，令 $X_j^{(A)}=X_j^{(B)}=1$，即可得到普通点积累加。

式~\eqref{eq:blockdot} 给出分层共享的算术边界；内部位宽、顺序与舍入仍由具体实现约束。



\subsection{通路复用与设计空间探索}
\label{sec:datapath_sharing}

表~\ref{tab:related_features} 汇总格式与组合能力。MX 列区分完整微缩放支持与仅元素运算支持；本文通过在共享 DP 内增配专用 Group 改变算力配比。

已有工作通过跨精度共享降低单元成本：TransDot 重组乘法器与移位器~\cite{transdot2026}，MPICC 将多种格式映射为统一内部表示~\cite{huangMPICCMultiplePrecisionInterCombined2025}，融合 FP8 点积设计复用累加和舍入后端~\cite{lutzFusedFP84Way2024}。DSE 研究进一步联合探索数值格式与 MAC 微架构~\cite{paretomac2025}，或共享尺度、权重和激活格式~\cite{zhao2026mxlike}；Timeloop 则以循环分块与存储层次分析阵列映射~\cite{timeloop2019}。本文固定目标格式集合，将 H/Q/E 峰值配比作为设计变量，考察实测单元成本如何通过面积与功耗预算影响部署规模与负载收益。

\subsection{算力配比的设计动机}
\label{sec:commercial_ratios}
\label{sec:ratio_motivation}

表~\ref{tab:vendor_ratios} 展示了从 $1:2:4$ 向 E 增强或 Q/E 同时增强的产品与研发路线。已有选择给出了配比变化的方向，但不能覆盖不同负载和硬件预算下的全部选择。本文进一步考察只增强 Q、提高 E 增配幅度以及扩大 Q/E 共同增配的组合，将 $1:4:4$、$1:2:12$、$1:2:16$、$1:6:12$ 和 $1:8:16$ 纳入探索，与上述四种配比共同构成九候选设计空间。



更高的低精度峰值需要额外乘法、对齐和规约资源，也会改变整核面积和逐模式功耗。因此，配比选择需要同时回答：单位成本能增加多少低精度吞吐，以及面积与功耗预算允许配置多少核心。本文以组合式微架构构造候选，以同工况测量建立成本，再将预算决定的绝对算力映射到 H/Q/E 负载组成，分析工业界已有选择和更大增配方案各自的优势区域。
